% Chapitre 10 : où déposer.
\begin{figure}[tbp]
\centering
\begin{tikzpicture}[x=1cm,y=1cm,
  case/.style={draw=black!40, fill=black!3, minimum width=4.3cm, minimum height=1.6cm,
    font=\sffamily\footnotesize, align=center, text width=4cm},
  tete/.style={font=\sffamily\scriptsize\bfseries, align=center, text=black!70}]
  \node[case] (a) at (0,0) {GitHub, Palomar\\[-1pt]\scriptsize\color{black!60} le code Lean, les
    données};
  \node[case] (b) at (4.5,0) {Zenodo\\[-1pt]\scriptsize\color{black!60} une version
    archivée, souvent avec un DOI};
  \node[case, draw=tonregle, fill=tonregle!7] (d) at (4.5,-1.8) {arXiv, HAL,
    ResearchGate\\[-1pt]\scriptsize\color{black!60} la prépublication qu'on cite};
  \node[case, fill=white, draw=black!25, text=black!50] (c) at (0,-1.8) {pas de voie
    fiable\\[-1pt]\scriptsize on y renvoie depuis la prépublication};
  \node[tete] at (0,1.15) {Code et données};
  \node[tete] at (4.5,1.15) {Prose};
  \node[tete, anchor=east, text width=2.5cm, align=right] at (-2.35,0)
    {non indexé par\\Google Scholar\\[1pt]\mdseries communiquer un savoir};
  \node[tete, anchor=east, text width=2.5cm, align=right, text=tonregle] at (-2.35,-1.8)
    {indexé par\\Google Scholar\\[1pt]\mdseries entrer dans le graphe des citations};
\end{tikzpicture}
\caption{Où déposer. La ligne du haut suffit pour qu'on trouve le travail, lecteurs et
modèles de langage compris ; seule la ligne du bas le fait entrer dans le graphe des
citations. On combine les deux.}
\label{fig:deposer}
\end{figure}
